\section*{第 \ref{ch:g7:central} 章——中心对称与平行四边形}

\begin{solution}{exo:g7:central:1}
\omterm{def:g7:central:def}{中心对称}的像：$M'$ 位于 $O$ 的\emph{左}侧四个方格、
\emph{下}侧一个方格处（两个方向都与原来相反）。关于竖直线的
轴对称像在左侧四个方格处，
但仍在\emph{上}侧一个方格处。
两者不同，它们互为上下方向的镜像。
\end{solution}

\begin{solution}{exo:g7:central:2}
$A'B' = AB$（\omterm{def:g7:central:def}{中心对称}保持长度），且
$(A'B') \parallel (AB)$（直线的像是与之\omterm{def:g6:lines:perp}{平行}的直线，
见\cref{prop:g7:central:preserved}）。
\end{solution}

\begin{solution}{exo:g7:central:3}
具有\omterm{def:g7:central:center}{对称中心}的计算器数字有：$0$、$2$、$5$、$8$（
$1$ 若画成普通竖条也算）。字母中，H、N、S、Z、O 有\omterm{def:g7:central:center}{对称中心}，A 和
T 则没有（它们有\omterm{def:g6:symmetry:axis}{对称轴}）。
\end{solution}

\begin{solution}{exo:g7:central:4}
对边相等：$MN = KL = 7$ cm，$NK = LM = 4$ cm。
对角相等：$\widehat M = \widehat K = 115^\circ$。
对角线在公共中点处互相平分，这一点即为 $O$：
\omterm{def:g6:lines:objects}{线段} $[LN]$ 的中点是 $O$。
\end{solution}

\begin{solution}{exo:g7:central:5}
作\omterm{def:g6:lines:objects}{线段} $[AC]$，使其长为 $8$ cm，标出其中点 $O$，
过 $O$ 作一条直线，夹角为 $60^\circ$，与 $(AC)$ 相交，并在其上取 $B$ 和 $D$，
各距 $2.5$ cm，分居 $O$ 两侧。依次连接 $A$、$B$、$C$、$D$。
由于对角线在各自中点处互相平分，故按定义 $ABCD$ 是
\omterm{def:g7:central:parallelogram}{平行四边形}。
\end{solution}

\begin{solution}{exo:g7:central:6}
邻角互补：与 $115^\circ$ 相邻的角为
$180 - 115 = 65^\circ$。对角相等，因此四个角分别为
$115^\circ$、$65^\circ$、$115^\circ$、$65^\circ$（和为
$360^\circ$）。
\end{solution}

\begin{solution}{exo:g7:central:7}
由 $A(1,1)$ 到 $O(2.5, 2.5)$，向右 $1.5$、向上 $1.5$；沿同样方向
继续，得 $A'(4, 4)$。由 $B(4,2)$ 到 $O$，向左 $1.5$、向上 $0.5$；
同样继续，得 $B'(1, 3)$。
\end{solution}

\begin{solution}{exo:g7:central:8}
不是。等腰梯形的两条腰相等，
但它不是\omterm{def:g7:central:parallelogram}{平行四边形}（两条\omterm{def:g6:lines:perp}{平行}边的长度不同）。
能够判定的是
\cref{thm:g7:central:recognize}中的条件 3：一组对边相等
\emph{且\omterm{def:g6:lines:perp}{平行}}。
\end{solution}

\begin{solution}{exo:g7:central:9}
$O$ 是 $[BC]$ 的中点（按取法），且由作图可知它也是
$[AA']$ 的中点。\omterm{def:g4:geometry:polygon}{四边形} $ABA'C$
的对角线 $[AA']$ 与 $[BC]$，有公共中点 $O$，依判定条件
1（即定义），无需额外条件即可判定它是\omterm{def:g7:central:parallelogram}{平行四边形}。
\end{solution}

\begin{solution}{exo:g7:central:10}
“对角相等 $\Rightarrow$ \omterm{def:g7:central:parallelogram}{平行四边形}”：\emph{正确}
（这是另一条判定条件；由角和
为 $360^\circ$，两组对角相等迫使邻角互补，
从而由
\cref{thm:g7:triangles:equalities}可知对边\omterm{def:g6:lines:perp}{平行}）。

“对角线\omterm{def:g6:lines:perp}{垂直} $\Rightarrow$ \omterm{def:g6:shapes:quadrilaterals}{菱形}”：对\omterm{def:g7:central:parallelogram}{平行四边形}而言\emph{正确}
（其对角线已在中点处互相平分）。

“有一个\omterm{def:g3:shapes:rightangle}{直角} $\Rightarrow$ 长方形”：\emph{正确}——
邻角互补，故四个角均为\omterm{def:g3:shapes:rightangle}{直角}。
\end{solution}

\begin{solution}{exo:g7:central:11}
无论\omterm{def:g4:geometry:polygon}{四边形}是怎样的——哪怕非常不规则——中点
$I$、$J$、$K$、$L$ 总构成一个\emph{\omterm{def:g7:central:parallelogram}{平行四边形}}。
证明借助中点定理，在\cref{ch:g8:midpoints}中
（见\cref{exo:g8:midpoints:9}）。
\end{solution}

\begin{solution}{pb:g7:central:1}
\textbf{1.} $M(3, 1) \to (-3, -1)$；$N(-2, 4) \to (2, -4)$；
$P(0, -5) \to (0, 5)$。规律：绕原点的\omterm{def:g7:central:def}{中心对称}
把 $(x, y)$ 变为 $(-x, -y)$——两个坐标都改变符号。

\textbf{2.} 从 $M(5, 3)$ 出发：位于 $C$ 右侧三个方格处的
点变为左侧三个方格，原来上方两个方格变为下方两个方格，像为
$(-1, -1)$。验算：把中心坐标的两倍减去该点坐标，
$(2 \times 2 - 5,\ 2 \times 1 - 3) = (-1, -1)$。

\textbf{3.} 绕卡片中心的\omterm{def:g7:central:def}{中心对称}把牌面上的圆点
两两互换。圆点数为\omterm{def:g3:numbers:evenodd}{奇数}时，必有一个圆点无人配对，它
只能与自身对应，而唯一与自身对应的点
就是中心本身。因此 3、5、
7\,\dots\ 的中间圆点都恰好位于正中央（看一张真卡片便知）。

\textbf{4.} 关于 $O_1(2,0)$：$A(1,1) \to (3,-1)$；关于
$O_2(5,0)$：$(3,-1) \to (7,1)$。起点 $A(1,1)$，终点
$(7,1)$：向右平移 $6$，而非翻转。对于
$B(0,3) \to (4,-3) \to (6,3)$，同样向右平移 $6$。两次
\omterm{def:g7:central:def}{中心对称}的复合相当于一个\emph{平移}。

\textbf{5.} 平移量为 $6$，而 $O_1 O_2 = 3$：平移
量恰为两中心距离的\emph{两倍}（沿
由 $O_1$ 到 $O_2$ 的方向）——正如两面\omterm{def:g6:lines:perp}{平行}镜子
平移其间距的两倍（见\cref{pb:g6:symmetry:1}）。

\textbf{6.} 绕 $O$ 的半转将\omterm{def:g7:central:parallelogram}{平行四边形}映到
自身（这正是“\omterm{def:g7:central:center}{对称中心}”对它的含义，
见\cref{thm:g7:central:properties}）。过 $O$ 的直线也被
映到自身，故它与边界相交的
两点彼此互换：它们关于 $O$ 对称。直线两侧的
两块\omterm{def:g7:central:parallelogram}{平行四边形}也互换，
因此它们全等（见\cref{prop:g7:central:preserved}）
——\omterm{def:g6:measure:area}{面积}相等。

\textbf{7.} 沿过\emph{两个中心}的直线剪开：
即蛋糕长方形的中心与空洞的中心。
由第 6 题用于蛋糕长方形，此剪线将蛋糕
平分；用于空洞长方形时（该直线同样经过其
中心），它也平分空洞。每一块都得到
蛋糕的\omterm{def:g3:division:half}{一半}减去空洞的\omterm{def:g3:division:half}{一半}：完全公平。

\textbf{8.} $D = (0 + 8 - 6,\ 0 + 5 - 1) = (2, 4)$（对角线
互相平分，故 $D$ 是 $B$ 的
关于中心的半转像）。中心为 $[AC]$ 的中点，
$O(4,\ 2.5)$。

\textbf{9.} 在\omterm{def:g7:central:parallelogram}{平行四边形} $ABA'C$ 中，对边
相等：$BA' = AC$。在 $ABA'$ 中，由三角形不等式
可得 $AA' < AB + BA' = AB + AC$。由于 $O$ 是
$[AA']$ 的中点，$AA' = 2\,AO$，故 $2\,AO < AB + AC$，即
$AO < \frac{AB + AC}{2}$。

\textbf{10.} 上界：$AO < \frac{5 + 7}{2} = 6$ 厘米。下界
：在 $ABA'$ 中，$AA' > BA' - AB = 7 - 5 = 2$，故
$AO > 1$ 厘米。自 $A$ 出发的中线严格介于 $1$ 与
$6$ 厘米。

\textbf{11.} \omterm{def:g6:lines:objects}{线段}：有，即其中点。直线：有——其上每一点
都是\omterm{def:g7:central:center}{对称中心}。\omterm{def:g6:shapes:triangles}{等边三角形}：没有。正方形：
有，即对角线的交点。圆：有，即其圆心。
S 与 Z：有，即它们的中点（把页面倒过来看：
两者看起来相同）。

\textbf{12.} 若存在\omterm{def:g7:central:center}{对称中心}，半转会使三个\omterm{def:g5:solids:def}{顶点}彼此配对；
三为\omterm{def:g3:numbers:evenodd}{奇数}，必有一个\omterm{def:g5:solids:def}{顶点} $A$ 为其自身的像，
从而 $A$ 必为半转的中心。另外两个\omterm{def:g5:solids:def}{顶点} $B$ 与 $C$
则关于 $A$ 对称——
于是 $A$ 为 $[BC]$ 的中点，三个\omterm{def:g5:solids:def}{顶点}
共线。三角形不可能三\omterm{def:g5:solids:def}{顶点}共线：
矛盾。故不存在\omterm{def:g7:central:center}{对称中心}。

\textbf{13.} 同样的奇偶性论证：保持\omterm{def:g4:geometry:polygon}{多边形}不变的半转
使其\omterm{def:g5:solids:def}{顶点}两两配对；\omterm{def:g5:solids:def}{顶点}数为\omterm{def:g3:numbers:evenodd}{奇数}时，
必有一个不动的\omterm{def:g5:solids:def}{顶点}，它须是中心，也是
另外两\omterm{def:g5:solids:def}{顶点}之间\omterm{def:g6:lines:objects}{线段}的中点——三个共线的\omterm{def:g5:solids:def}{顶点}，
这在\omterm{def:g4:geometry:polygon}{多边形}中不可能\,\dots\ 故
任何\omterm{def:g5:solids:def}{顶点}数为\omterm{def:g3:numbers:evenodd}{奇数}的\omterm{def:g4:geometry:polygon}{多边形}都没有\omterm{def:g7:central:center}{对称中心}。

\textbf{14.} 先绕 $O_1$ 作半转，再绕
$O_2$ 作半转，两次都把图形映回自身——因此二者的
复合也如此。由第 5 题，该复合是一个平移，
平移量为 $O_1 O_2$ 的两倍，它不为零，
因两中心不同。能放在一张纸上的图形，
不可能与自身平移某一固定非零量后的图形相等
（多次平移后必将离开纸面）。故有界图形
至多只有一个\omterm{def:g7:central:center}{对称中心}。

\textbf{15.} 在无限长的条带上，绕位于左脚印与其后
右脚印之间的中点的半转将图案映回自身；
再向前一整步的点同样如此：
两个不同的\omterm{def:g7:central:center}{对称中心}。两次半转的复合
得到平移量为半步两倍的平移——恰好
是明显将无限长的花边映到
自身的一步平移，正如第 14 题所预言。无界图案遵循
另一套规则：这正是壁纸的数学。
\end{solution}
